\section{相关研究现状}

\todo{1.核实文献中的方法论总结是否正确妥当；2.改为一段总结一次现有研究局限性的写法，并在最后一段中使用文献[x-y]的方式引用所有相关工作，而不是列举一篇批判一篇}

针对向LLM Agent工具调用链的间接提示注入攻击威胁，现有防御工作大致沿着“提示工程与输入净化—训练与安全对齐—Agent系统架构设计”的技术路径演进，以下从这三个层面分别梳理现有研究工作的脉络及其局限。

\subsection{基于提示工程与输入净化的防御方法}

早期研究试图仅通过提示文本的组织方式或简单的净化规则抵御间接提示注入，其共同特点是既不改变模型参数，也不改变Agent的执行架构，而是依赖提示层面的结构调整引导模型区分指令与数据。Schulhoff\cite{schulhoff2024sandwich}提出的Prompt Sandwiching（提示夹心法）在不可信内容前后重复插入用户原始指令，试图借助模型对近端上下文的注意力偏好，使其在读取外部内容后仍能回忆并优先执行最初的任务目标；该方法未对不可信内容本身做任何标记或隔离，防御效果完全取决于重复指令能否在生成过程中持续压制注入指令的影响，对措辞隐蔽或与任务高度相关的注入内容缺乏针对性。Hines等\cite{hines2024defending}提出的Spotlighting（数据分隔符法）则通过对工具返回的不可信内容施加特殊分隔符标记或字符编码变换（如虚假语言编码、编码转写），使模型有机会将其显式识别为“数据”而非“指令”，从而降低对其中指令性内容的遵从概率，但该方法仍要求模型具备足够的指令遵从与格式识别能力才能生效，一旦攻击者在注入载荷中模仿或破坏分隔符格式，该标记信号即可能失效。Zhang等\cite{zhang2025asb}提出的Agent Security Bench基准进一步将分隔符、提示夹心法、指令预防（Instructional Prevention）、复述（Paraphrasing）与乱序（Shuffle）等同类简单方法系统化为统一的基线家族，并纳入同一评测协议中比较其防御效果，其结果表明这类方法在受控场景下可降低部分攻击成功率，但普遍存在效果不稳定、随攻击话术变化而大幅波动的问题。Shi等\cite{shi2025promptarmor}提出的PromptArmor则在数据净化侧引入了更工程化的检测流程，将输入内容交由专用的外部护栏LLM做净化处理，但其检测仍是基于语言模型的语义判断，完全依赖于护栏LLM的识别能力与鲁棒性，难以处理外部护栏LLM不可信或检测模式被对抗性绕过的风险。

上述基于提示工程与输入净化的防御方法的共同局限在于：防御效果完全依赖于LLM自身对分隔符、重复指令或结构化标记的遵从与识别能力，本质上是一种“软约束”而非强制性的执行控制，无法提供确定性的安全边界，难以为Agent工具调用链提供可验证的安全保证。

\subsection{基于训练的防御方法}

为弥补提示工程类方法的不足，已有一批研究转向通过训练或微调的方式，使大模型或专用检测器具备更稳定的指令边界识别与抗注入能力。在大模型层的安全对齐方面，Wallace等\cite{wallace2024instruction}提出的Instruction Hierarchy显式定义系统消息、用户输入与第三方工具/数据输出之间的优先级，并提出“上下文综合”与“上下文忽略”两种自动数据生成方法，训练模型在指令发生冲突时优先遵循系统级特权指令、忽略或驳回与其不一致的低特权内容；该方法显著提升了模型对指令层级的敏感度，但其约束仍建立在模型参数化的概率生成之上，无法保证在所有输入分布下均严格遵守优先级规则，对训练数据未覆盖的攻击话术泛化能力有限。在输入结构化方面，Chen等\cite{chen2025struq}提出的StruQ在LLM系统前构建前端，通过特殊标记将指令与数据分离到不同的结构化通道，并微调专门的模型识别这些标记，以降低数据被误解析为指令的概率；但该机制依赖于前端标记未被绕过的假设，一旦攻击者构造的内容能够穿透或伪造分隔标记，防御即可能被突破，且需要针对系统输入分隔前端对基座模型进行专门的微调以适应其数据处理结构，在不同应用场景下的部署和迁移成本较高。Chen等提出的SecAlign\cite{chen2025secalign}和Meta SecAlign\cite{chen2025metasecalign}将提示注入防御进一步形式化为偏好优化的对齐训练问题，采用直接偏好优化范式训练模型在面对含注入指令的输入时偏好安全输出、抑制不安全响应，在保持常规任务能力的同时提升了对提示注入的鲁棒性，但两者的安全性均以对齐训练所覆盖的攻击分布为边界，针对未见攻击的泛化能力有限。Mou等\cite{mou2026toolsafe}提出的ToolSafe框架引入基于多任务防御反馈强化学习微调的步骤级主动安全护栏模型，在每步工具调用前进行安全检查，但其安全判断仍以提示形式注入上下文中，反馈机制仍然是松耦合的，无法强制约束Agent拒绝危险工具调用。而在提示注入文本检测模型方面，ProtectAI\cite{protectai2024}提出的基于DeBERTa-v3-base微调的提示注入分类器作为业界广泛使用的轻量级检测基线，以较低的推理开销对输入文本进行二分类判别；Li等\cite{li2025piguard}提出的PIGuard针对已有提示注入护栏普遍存在的“过度防御”问题（即将大量良性输入误判为攻击）进行专门优化，在维持较高攻击检出率的同时降低了良性任务的误报率，但这些方法作为独立于Agent执行流程的外挂式检测器，其判别范围局限于孤立的文本片段，且主要局限于“用户输入中是否包含提示注入攻击词”这一检测场景，而非针对LLM Agent执行场景中跨工具调用步骤的语义传播与动作风险设计。Liu等\cite{liu2025datasentinel}提出的DataSentinel则提出了一种博弈论意义下的极小极大（minimax）检测训练框架，将检测器训练建模为攻击者与检测器之间的对抗博弈，以提升检测器对适应性攻击的鲁棒性；但该方法的检测粒度仍停留在单次输入的静态判别层面，而非LLM Agent场景普遍存在的动态、多轮的工具调用链。整体来看，基于训练的防御方法相较提示工程类方法提供了更稳定的统计意义上的安全收益，但其安全性本质上仍是模型或检测器在训练数据分布内的经验性泛化，对训练阶段未见过的攻击或自适应攻击的防御性能尚未得到验证，且训练与更新成本远高于提示工程类方法，在部署级系统中难以随攻击方法的快速演化及时调整。

\subsection{基于Agent系统架构设计的防御方法}
\todo{核实未分类的文献内容，如果发现有错进行修正，如果发现是新分类可以留下，如果发现不能紧扣研究主题的的则删去保持文章内容简洁}

传统系统安全领域中，污点分析（Taint Analysis）与溯源图（Provenance Graph）长期被用于追踪不可信数据在程序执行中的传播路径、关联系统事件的因果依赖\cite{zipperle2022provenance}，为具有相似因果结构的LLM Agent工具调用链安全研究提供了方法论借鉴。然而，Agent场景下数据传播由概率性的自然语言推理而非语法依赖决定，传统方法难以直接迁移，因此近年研究开始关注运行时监控、执行溯源与因果归因以及架构级信息流控制三个方向，探索面向Agent工具调用链的专门化防御机制。

\textbf{运行时监控与执行隔离。}这类工作的防御时间节点在工具调用的实际执行环节。Shi等\cite{shi2025progent}提出的Progent权限控制框架，运行时拦截器使用SMT求解方式在每次工具调用前判定工具权限是否单调衰减，以控制“权限升级”攻击，但其主要防御“工具权限集合不会未经批准扩大”，而难以防御“在合法权限内完成”的攻击，如偏好操纵攻击。Li等\cite{li2025drift}提出的DRIFT采用基于LLM的“安全规划器生成预期工具调用轨迹—动态验证器逐步比对—注入隔离器清洗外部注入指令”的三级架构，在运行时动态审批偏离预期规划的工具调用，但其防御有效性高度依赖安全规划器对任务轨迹预判的准确性，若规划器本身受到不可信输入影响，则可能导致防御被绕过，且三级验证带来了显著的计算开销与任务延时。Zhu等\cite{zhu2025melon}提出的MELON采用掩码后重执行的思路，通过在主Agent执行过程中引入一个额外的影子执行分支，将用户输入掩码替换为中性的无关指令，观察掩码前后Agent工具调用决策是否受到外部注入影响，来判断是否受到注入攻击，但步骤级掩码重执行机制同样引入了较高的计算开销与延时。Wang等\cite{wang2025agentspec}提出的AgentSpec提出面向运行时约束的领域特定语言，允许开发者将安全策略转化为形式化执行规则，在代码执行、具身Agent等高风险场景中实现了低开销下的高覆盖率拦截；但其策略依赖预先人工编写，面对未见场景难以自动扩展。Wu等\cite{wu2025isolategpt}提出的IsolateGPT将隔离机制下沉至执行环境层，为每个工具分配独立沙箱上下文以切断工具间的状态共享与权限继承，但该方法引入了额外的性能开销与部署成本，且防御边界受限于沙箱实现本身，若沙箱实现不善，则依旧有逃逸和绕过的风险\todo{核实并准确评估针对IsolateGPT的说法}。Wen等\cite{wen2026agentsys}提出的AgentSys通过显式的分层记忆管理，在隔离的子任务执行单元内对命令类工具调用触发验证器，检测到风险后对当前观察内容进行净化并在原意图下有界重启，但其验证机制仅覆盖隔离子任务内部的调用路径，顶层Agent发起的工具调用未经过同一验证器，顶层Agent仍然存在被注入的风险\todo{此说法存疑，重新核实}。Wang等\cite{wang2025cacheprune}提出的CachePrune从模型内部表示层面出发，通过在推理时监控并对KV缓存中与注入指令相关的注意力模式进行编辑，引导模型不遵循被识别出的注入内容，该方法无需重新训练，但其防御粒度与具体模型的缓存结构和内部表示绑定，针对实际部署场景下跨模型架构或黑盒API场景下的可迁移性有限。而且，总体而言，运行时监控类方法仍以单步“阻断或放行”的二值决策为主要输出，在损害任务效用的同时，也难以支撑精细化的防御响应或事后审计。

\textbf{执行溯源与因果归因。}这类工作试图为Agent的运行轨迹构建可查询的溯源结构，进而定位风险的来源与传播路径。An等\cite{an2025ipiguard}提出的IPIGuard通过建模工具间的正常依赖关系，以异常依赖边识别注入攻击导致的非预期工具调用，证明了工具调用依赖结构可作为识别外部不可信信息注入的有效信号；但其工具依赖图需要在任务执行前预先规划，对调用顺序难以预先确定的动态复杂任务适用性不足。Wang等\cite{wang2025agentarmor}提出的AgentArmor将Agent运行时轨迹抽象为程序依赖图，并引入安全类型系统进行细粒度信息流控制与策略检查，在多个基准上取得了较高的攻击拦截率，但其复杂的图构建过程和逐调用的LLM依赖边推理带来了显著的额外计算开销，且用于推断依赖关系的LLM本身也可能遭受二次注入攻击。Wang等\cite{wang2026authgraph}提出的AuthGraph构建授权图与溯源图相结合的双图架构，以参数策略约束关键参数的允许来源、并将动态重规划限制在白名单与程序检查范围内；但其授权图本身在隔离的干净上下文中生成，并不能独立于生成过程构成授权真值。Cai等\cite{cai2026ghost}提出的NeuroTaint将关注点从语法依赖的信息流传播转向基于语义转换与因果推理的传播，设计动态上下文溯源图以在Agent记忆边界和会话重启后维持污点谱系的持久化，并引入汇点驱动的因果分析器通过反事实推理捕捉隐式控制流影响；但该方法定位为离线的事后审计工具，无法在Agent执行时实时阻断或净化危险行为。Sequeira等\cite{sequeira2026agentsentry}提出的Agent-Sentry从Agent历史良性轨迹中学习正常行为边界，构建“结构分类器初判—白名单参数来源过滤—LLM残余仲裁”的三层运行时门控架构，将大部分判定固化在确定性结构检查中以降低LLM调用频率；但其参数来源证据由同一Agent自行上报，来源可信度本身仍缺乏独立验证\todo{继续详细核实这个说法}。Zhang等\cite{zhang2026agentsentry}提出的AgentSentry，其针对工具返回边界设计了受控的反事实“空跑”（dry-run）机制，通过因果诊断判断工具返回内容是否驱动了后续风险行为，并以因果门控的方式对上下文进行净化后继续任务\todo{完善说法}。Kim等\cite{kim2026causalarmor}提出的CausalArmor对特权决策执行留一法（Leave-One-Out）意义下的归因分析，仅在观测到“主导权转移”（dominance shift）时才触发针对性净化，以较低成本实现选择性归因；但其归因结论本质上是移除特定片段后的操作性反事实影响，并不构成严格意义上的一般因果证明\todo{核实此说法}。Yan等\cite{yan2026propguard}提出的PropGuard采用双视图时空图恢复有害传播子图，对定位到的来源节点执行抑制、净化或重新生成，并按依赖顺序对下游交互进行重放；但该方法主要停留在图结构层面的传播定位与重放验证，缺少对切断效果的运行时物化回执，难以确认危险路径在真实执行层面已被永久阻断\todo{核实此说法}。上述工作表明，执行溯源与因果归因方法能够刻画注入攻击在Agent行为轨迹中难以完全伪装的结构指纹，相比纯语义检测具有更强的鲁棒性，但现有工作普遍存在“检测有余、闭环不足”的问题，主要聚焦于风险定位或事后审计，缺乏与阻断、恢复动作的一体化验证。

\textbf{架构级信息流控制。}针对运行时方法在安全保证上的概率性局限，一系列工作从系统架构层面重新设计Agent的执行模型，试图以确定性机制取代概率性过滤。Debenedetti等\cite{debenedetti2025camel}提出的CaMeL提出规划与执行分离的双LLM架构，特权规划器仅读取用户提示生成数据流图，被隔离的执行器严格按图调用工具，从架构层面阻断了不可信数据对控制流的劫持，但其代价显著高昂，任务效用损失、额外LLM的计算成本和单次防御延时较高，不适合用户任务场景。Costa等\cite{costa2025fides}提出的FIDES将信息流控制方法应用于Agent，为内容附带完整性/机密性标签并强制传播，在敏感工具执行前对不可信数据进行隔离提取与策略门控；但该方法只能阻止显式信息流，无法阻止由控制流决策引起的隐式流，且其规划器完全依赖LLM进行污点数据改写提取，计算开销较高、可扩展性不佳。Kim等\cite{kim2025promptflow}提出的Prompt Flow Integrity借鉴系统安全中的最小权限原则，将Agent拆分为处理可信与不可信数据的双组件，通过跨组件通信监控与显式/隐式流检查在不可信数据流向敏感汇聚点时触发告警或阻断；但该方法需要在插件层面手动标注数据源并管理代理与流检查策略，任务效用受限、策略规模化与跨服务追踪困难。Wu等\cite{wu2024fsecure}提出的f-secure LLM系统从信息流控制的视角出发，将LLM系统解耦为具备上下文感知能力的处理管线，并依据信息流控制原则动态生成结构化的执行方案，试图从系统级别限制不可信数据对控制流的影响；但该方案同样依赖于对数据来源与信任等级的准确标注，标注误差会直接传导至执行方案的安全性。这类方法整体代表了Agent工具安全从概率性过滤走向确定性约束的趋势，但普遍面临部署复杂度高、任务效用折损显著、跨框架集成困难等工程挑战，且多数方法仍以阻断攻击为核心目标，缺少对风险传导路径、归因分析与残余风险的结构化解释。

\todo{核实以下这些方法的分类、内容描述及局限性评述是否准确}
Zhang\cite{zhang2026agentlibos}提出的Agent libOS借鉴操作系统理念，提出“进程身份+能力+运行时基元+审计”的权威边界，使资源权威仅在显式、可审计的能力调用原语内改变；但其进程与能力语义需在宿主操作系统与语言运行时之上构建新的抽象层，部署与适配成本较高\todo{不要使用“+”的连接方式声明方法贡献，在无法说明资源权威、可审计、原语等名词是什么的情况下不要滥用，可以转写说法}。Fan等\cite{fan2026granularity}提出的PACT从工具模式合成参数角色契约，跨重规划追踪值的来源，在工具调用前逐参数核验信任度、禁止来源与义务约束，为可信过程提供了有限的凭证化豁免机制；但其自动角色与来源推断的准确率仍存在明显误差，且该机制作用于参数粒度，未覆盖任务级别的整体调用授权\todo{未解释的情况下不要使用“跨重规划”“可信过程”这样的名词，考虑转写}。Santos-Grueiro\cite{santosgrueiro2026temporary}提出的Commit-Time Authorization在提交时刻绑定授权见证、派生状态与持久效果，核验新鲜度、因果优先级与效果绑定后支持刷新、重新绑定、重规划或拒绝；但该机制本身只保证提交状态与见证记录之间的一致性，并不单独判定特定内容是否构成语义上已获授权的间接提示注入。Hossain等\cite{hossain2026nexus}提出的NEXUS在参数级检查后输出允许、阻断、确认或修订四类动作；但其压力测试表明该方法在应当输出“修订”的样本中未能正确预测出任何一例，多数情况回退为“确认”，暴露出细粒度动作校准上的明显不足。Ma等\cite{ma2026secureclaw}提出的SecureClaw通过句柄限定与“预览—提交”两阶段绑定机制，在提交前由可信执行器重新核验请求的新鲜度、重放与确认状态，拒绝后仅返回粗粒度动作类别与安全提示；但被阻断的提交无法清除已经进入规划器上下文中的明文内容，遗留了信息层面的残余风险。Sun等\cite{sun2026triad}提出的TRIAD让训练后的护栏模型在每个规划步骤输出继续、更新或拒绝三类判定并附带结构化反馈，驱动Agent在执行前对计划进行修订；但其护栏基于特定规模的训练模型、且在模拟工具与离线评测环境下验证，额外推理时延与对未见攻击的泛化能力仍待进一步检验。

\todo{二次核实本段说法，尤其是“残余风险”“泛化能力”等表述是否准确，不得有不符合事实的overclaim与绝对化表述}
综合上述工作可以看出，现有面向Agent工具调用链的提示注入攻击防御沿着提示工程、训练对齐与系统架构三条路径不断发展，呈现出两端分布的特征：提示工程与训练类方法部署成本较低，但其安全性建立在模型自身概率性遵从能力之上，在未见攻击上防御效果显著下降，难以提供鲁棒的安全性；而架构级信息流控制、可编程权限引擎与执行隔离类方法虽能提供更强的确定性约束，却普遍伴随较高的运行时开销、复杂的策略编写与维护负担，部分方法的授权判定还依赖对来源真实性或结构化账本正确性的假设，安全性与任务效用之间的矛盾尚未得到很好的调和。另一方面，现有的执行溯源、因果归因与包含性评测工作大多分别聚焦于风险检测、归因定位或状态恢复中的某一环节，尚未在同一评测框架下系统评估授权判定、跨步骤攻击链遏制与执行恢复三者的联合效果，也尚未充分探讨在调用被拒绝后向规划器披露何种粒度的反馈信息，才能在保持安全性的同时尽量维持任务效用这一问题。上述两方面的不足，为本研究分别从任务契约驱动的工具调用前防御与运行轨迹图驱动的在线检测与恢复两个研究点展开工作提供了动机。